\documentclass[11pt]{article}
\usepackage[margin=1in]{geometry}
\usepackage{amsmath,amssymb,amsthm,mathtools,microtype,hyperref}
\newtheorem{theorem}{Theorem}
\newtheorem{lemma}[theorem]{Lemma}
\newtheorem{corollary}[theorem]{Corollary}
\newtheorem{remark}[theorem]{Remark}
\newcommand{\E}{\mathbb E}
\newcommand{\1}{\mathbf 1}
\newcommand{\cE}{\mathcal E}
\DeclareMathOperator{\sgn}{sgn}
\title{Homogeneous parity approximants have a strict local repair cost}
\author{Dmitry Vaintrob \and Kaarel H\"anni \and OpenAI Codex}
\date{October 2026}
\begin{document}
\maketitle
\begin{abstract}
We study a discrete cusp-kernel model of local repairs to a globally
margin-separated prediction field. Repairs may depend on every active
coordinate and may change their signs arbitrarily many times before
returning to the background signs. We prove that a homogeneous polynomial
of sufficiently small degree, with parity error at most $0.34$, has mean
repair cost strictly above the constant-background benchmark. Amplitudes
need not be bounded. The proof uses an amplitude transport and
an inverse-amplitude normalization that remains valid for arbitrarily
small transport rows. A more general theorem identifies the precise
one-sided adjacency deficit that controls the cost. A weighted extension
also covers certain affine homogeneous heads after the exact odd
augmentation of the neural kernel. These are results
about a fixed-background local model; they do not establish the desired
sample lower bound for general adaptive neural-prior posteriors.
\end{abstract}

\section{Model and result}
Let $\Omega_k=\{-1,1\}^k$ have uniform probability measure, let
$f(x)=\prod_{i=1}^k x_i$, and write $d_H$ for Hamming distance. Fix
$0<\beta\le1/2$. The local index set is
$\mathcal X=\Omega_k\times\mathbb Z_{\ge0}$. Its second coordinate is a
discrete nuisance radius. For finitely supported real coefficients
$c_z$ satisfying $\sum_zc_z=0$, put
\begin{equation}\label{eq:variance}
 V_\beta(c)=-\sum_{(x,h),(y,j)}c_{x,h}c_{y,j}
                    \bigl(d_H(x,y)+h+j\bigr)^\beta.
\end{equation}
This quadratic form is nonnegative. Indeed, the identity
\[
 s^\beta=\frac{\beta}{\Gamma(1-\beta)}
       \int_0^\infty(1-e^{-st})t^{-1-\beta}\,dt
\]
reduces the assertion to positivity of the kernels
$e^{-td_H(x,y)}e^{-th}e^{-tj}$. The first factor is a product of
positive semidefinite two-point kernels.

For a real correction $b$ on $\mathcal X$, define its squared energy by
\begin{equation}\label{eq:energy}
 \cE_\beta(b)=\sup_{\substack{c\text{ finitely supported},\ \sum c=0\\
                              V_\beta(c)>0}}
                  \frac{\left(\sum_zc_zb(z)\right)^2}{V_\beta(c)}.
\end{equation}
If a zero-variance contrast has nonzero value, the energy is defined as
$+\infty$. Constants have zero energy. This is the Hilbert seminorm
associated with the conditionally positive kernel in
\eqref{eq:variance}; in particular, every contrast obeys
\begin{equation}\label{eq:cs}
 \left|\sum_zc_zb(z)\right|^2\le\cE_\beta(b)V_\beta(c).
\end{equation}

Fix a background $P:\Omega_k\to\mathbb R$ with $|P|\ge1$. Set
$g=\sgn P$, $B=|P|$, $E=\{g\ne f\}$, $C=E^c$, and $R=\Pr(E)$.
For each $u\in E$, an admissible repair is a correction $b_u$ for which
\begin{enumerate}
\item $Z_u(x,h)=P(x)+b_u(x,h)$ satisfies $|Z_u(x,h)|\ge1$ everywhere;
\item $\sgn Z_u(u,0)=f(u)$;
\item for every $x$, $\sgn Z_u(x,h)=g(x)$ for all sufficiently large $h$.
\end{enumerate}
No condition is imposed on intermediate signs. Write
$e_u=\cE_\beta(b_u)$ and
\[
 v_0=2-2^\beta,\qquad C_0=\frac4{v_0},\qquad T=\frac{11}{3},
 \qquad V=2^\beta(2-2^\beta)+\frac{2^{2\beta}-2^\beta}{24}.
\]
All averages below use the original uniform measure, so
$\E[\1_E e_u]$ includes the error probability.

For a correct vertex $v$, define
\begin{equation}\label{eq:deficit}
 H_v=\frac1k\sum_{u\sim v,\ u\in E}B_u,
 \qquad \delta_v=\left(1-\frac{H_v}{B_v}\right)_+,
 \qquad D_C=\E[\1_C\delta_v].
\end{equation}

Define also the error-to-correct edge mass and amplitude-ratio sum
\[
 F=\E\left[\1_E\frac{\#\{v\sim u:v\in C\}}k\right],
 \qquad M=\E[\1_C H_v/B_v].
\]

\begin{theorem}[Signed surplus and one-sided deficit]\label{thm:deficit}
For $k\ge2400$, every background and every family of admissible repairs
as above satisfy
\begin{equation}\label{eq:main-deficit}
 \E[\1_Ee_u]\ge2T(2R+M-F)-T^2VR
                  \ge2T(1-D_C)-T^2VR.
\end{equation}
There is no polynomial-degree or amplitude hypothesis in this statement.
\end{theorem}

\begin{corollary}[Homogeneous low degree]\label{cor:homogeneous}
Suppose $P$ is supported on a single Walsh level $r$, with
$r/k\le1/5$ and $|P|\ge1$. If $k\ge2400$ and $R\le17/50$, then
\begin{equation}\label{eq:homogeneous}
 \E[\1_Ee_u]\ge\frac{C_0}{2}+\frac25.
\end{equation}
The constants hold uniformly for $0<\beta\le1/2$.
\end{corollary}

The benchmark $C_0/2$ is exact for a constant background. To see this,
take $P=-1$, and at each wrong center let $s=d_H(u,x)+h$ and
\[
 r_1(s)=(s+1)^\beta-s^\beta,\qquad
 b_u(x,h)=2-\frac2{v_0}+\frac2{v_0}r_1(s).
\]
At the center $b_u=2$, while for integer $s\ge1$ monotonicity of
$r_1$ gives $b_u\le0$. Thus all required margins and eventual signs
hold. The function $r_1$ is the representer of the contrast between
$(u,0)$ and $(u,1)$ and has energy $v_0$. Hence this repair has energy
$C_0$. The constant background errs on half the cube.

\section{A dual inequality allowing all sign changes}
We first establish the geometric estimate used in the transport proof.

\begin{lemma}[Bounded neighbor dual]\label{lem:dual}
Fix a wrong center $u$. Let $w_v\ge0$ be supported on correct neighbors
of $u$, with $a=\sum_vw_v\le1$ and $s=\sum_vw_v^2\le1/12$. Put
\[
 L_w=\sum_vw_v(2+B_u-B_v),\qquad
 F_w=2^\beta(2-2^\beta)a^2+(2^{2\beta}-2^\beta)s.
\]
Every admissible repair satisfies $e_u\ge C_0$ and
\begin{equation}\label{eq:dual}
 e_u\ge2tL_w-t^2F_w\qquad(0\le t\le11/3).
\end{equation}
\end{lemma}
\begin{proof}
At a fixed active point, the contrast between radii $h$ and $h+1$ has
variance
\[
 v_h=2(2h+1)^\beta-(2h)^\beta-(2h+2)^\beta
 =\beta(1-\beta)\int_0^1\!\int_0^1
                 (2h+s+t)^{\beta-2}\,ds\,dt.
\]
This decreases with $h$. An adjacent-radius sign change has magnitude
at least $2$, so it costs at least $4/v_h$. The corrected center must
eventually return to its background sign, proving $e_u\ge4/v_0=C_0$.
For $h\ge1$,
\begin{equation}\label{eq:late}
 v_h\le\tfrac14(2h)^{-3/2},\qquad
 v_h\le v_1\le2\sqrt3-\sqrt2-2<\tfrac1{20}.
\end{equation}
The second bound uses that $v_1$ increases with $\beta$ on $(0,1/2]$: in
its integral representation both $\beta(1-\beta)$ and $(2+s+t)^{\beta-2}$
increase with $\beta$.

Correct neighbors of a wrong parity center have the same background
sign as the center. If their signs have returned to that background
sign by radius $H$, the contrast with coefficient $a$ at $(u,0)$ and
coefficients $-w_v$ at $(v,H)$ has value of magnitude at least
$(L_w)_+$. Direct substitution in \eqref{eq:variance} gives its variance
\begin{equation}\label{eq:FH}
 F_H=v_0a^2(H+1)^\beta+s\bigl((2H+2)^\beta-(2H)^\beta\bigr).
\end{equation}
Thus $(L_w)_+^2\le e_uF_H$.

If $e_u<80$, \eqref{eq:late} excludes every sign change at radius
$h\ge1$. Eventual anchoring therefore fixes all signs already at
$H=1$. Since $F_1=F_w$, completing the square proves
\eqref{eq:dual}, in this case for every $t\ge0$.

If $e_u\ge80$, choose
\[
 H=\left\lfloor\tfrac12(e_u/16)^{2/3}\right\rfloor+1.
\]
The first bound in \eqref{eq:late} excludes every sign change at or
beyond $H$. Moreover, $v_0 2^\beta\le1$, $\beta\le1/2$, and
$s\le1/12$ give
\[
 F_H\le\sqrt{(H+1)/2}+\frac1{12}
 \le\sqrt{1+\tfrac14(e_u/16)^{2/3}}+\frac1{12}
 \le\frac{e_u}{4(11/3)^2}.
\]
For the last inequality, the ratio of the left side to $e_u$ decreases
with $e_u$. At $e_u=80$, its numerator is less than $17/12$, whereas
$80/[4(11/3)^2]=180/121>17/12$.
Consequently $(L_w)_+\le e_u/[2(11/3)]$. This implies
$e_u\ge2tL_w\ge2tL_w-t^2F_w$ for the asserted range of $t$.
\end{proof}

\section{Amplitude transport with no discarded rows}
\begin{proof}[Proof of Theorem~\ref{thm:deficit}]
On an error-to-correct neighbor edge define
\begin{equation}\label{eq:flow}
 \mu_{uv}=\frac{B_u}{k}.
\end{equation}
This raw transport gives
\[
 \sum_v\mu_{uv}\le B_u,\qquad
 \sum_u\mu_{uv}=H_v,\qquad
 0\le\mu_{uv}\le B_u/k.
\]
Set
\[
 \theta_u=\frac{\sum_v\mu_{uv}}{B_u}\le1,\quad
 w^0_{uv}=\frac{\mu_{uv}}{B_uB_v},\quad
 a_u=\sum_vw^0_{uv},\quad J_u=B_ua_u-\theta_u.
\]
Here $\theta_u$ is exactly the fraction of correct neighbors of $u$.
Since $B_v\ge1$, we have $a_u\le1$ and $w^0_{uv}\le1/k$.
Zero rows are allowed and have $a_u=\theta_u=J_u=0$.

Let $\eta=1/100$, $s_u=\max(a_u,\eta)$, and apply
Lemma~\ref{lem:dual} to $w_{uv}=w^0_{uv}/s_u$. Its row sum is at most
$1$, and
\[
 \sum_vw_{uv}^2\le\frac{a_u}{ks_u^2}
                  \le\frac1{k\eta}\le\frac1{24}.
\]
Its linear contrast is $(J_u+2a_u)/s_u$ and its variance is at most
$V$. Taking $t=T$ yields
\[
 J_u\le s_u A_u-2a_u,\qquad
 A_u=\frac{e_u+T^2V}{2T}.
\]
The baseline bound $e_u\ge C_0$, together with $C_0\ge4$ and
$V\ge2\sqrt2-2$, implies
\[
 A_u\ge\frac{4+(121/9)(2\sqrt2-2)}{22/3}>\frac{200}{99}.
\]
In particular $A_u-2\ge0$ and $\eta A_u\le A_u-2$. If
$a_u\ge\eta$, then $J_u\le a_u(A_u-2)\le A_u-2$; if
$a_u<\eta$, then $J_u\le\eta A_u-2a_u\le A_u-2$.
Thus every row, including a zero or arbitrarily small row, obeys
\begin{equation}\label{eq:pointwise}
 J_u\le\frac{e_u+T^2V}{2T}-2.
\end{equation}

The inverse-amplitude weights now turn the transport into a count
identity:
\[
 \E[\1_E B_ua_u]=M,\qquad \E[\1_E\theta_u]=F.
\]
Thus $\E[\1_EJ_u]=M-F$. Averaging \eqref{eq:pointwise} proves
the first bound in \eqref{eq:main-deficit}. The second follows from
$M\ge1-R-D_C$ and $F\le R$. In particular, surplus incoming amplitude
at correct vertices has not been discarded.
\end{proof}

\section{Homogeneity and the explicit strict gap}
Let $A$ be normalized cube adjacency. A homogeneous Walsh polynomial
of degree $r$ satisfies $AP=(1-2r/k)P$. Multiplication by full parity
reverses every edge, so $q=fP$ satisfies
$Aq=-(1-2r/k)q$. At a correct vertex $v$, therefore,
\[
 H_v=(1-2r/k)B_v+\frac1k\sum_{w\sim v,\ w\in C}B_w
       \ge(1-2r/k)B_v.
\]
This already gives $D_C\le(2r/k)(1-R)$, but retaining the
correct--correct edges gives a substantially stronger bound. Pairing the
two orientations of every such edge and using $a/b+b/a\ge2$ yields
\[
 M=(1-2r/k)(1-R)+\E\left[\1_C\frac{A_C B}{B}\right]
 \ge (2-2r/k)(1-R)-F.
\]
Since $F\le R$, Theorem~\ref{thm:deficit} gives
\begin{equation}\label{eq:alpha}
 \E[\1_Ee_u]\ge4T(1-r/k)(1-R)-T^2VR.
\end{equation}

For completeness, the numerical gap in Corollary~\ref{cor:homogeneous}
can be checked algebraically. The right side of \eqref{eq:alpha}
decreases with both $r/k$ and $R$. At $r/k=1/5$ and $R=17/50$, its
first term is $968/125$, exceeding $2T=22/3$ by $154/375$.
Put $v=2-2^\beta\in[2-\sqrt2,1)$. Multiplying the excess of
$2T-T^2V(17/50)$ above $C_0/2$ by $v$ gives
\[
 F_0(v)=\frac{47311}{10800}v^3-\frac{2057}{240}v^2
                         +\frac{37543}{5400}v-2.
\]
The derivative has positive leading coefficient and discriminant
$-696300671/9720000<0$, so $F_0$ is increasing. In particular
\[
 F_0(v)\ge F_0(2-\sqrt2)
 =\frac{86567}{1800}-\frac{18359\sqrt2}{540}
 >\frac{121}{10000}.
\]
Because $v\le1$, the unscaled gap is at least this large. Therefore the
right side of \eqref{eq:alpha} exceeds $C_0/2$ by more than
$121/10000+154/375>2/5$, proving the corollary.

For a mixed-degree head, the same calculation identifies a sufficient
relative generator condition. With
$\Delta=\sum_i(I-\mathrm{flip}_i)/2$, at correct vertices
\[
 H_v=B_v-\frac2k\sgn(P_v)\Delta P(v)
          +\frac1k\sum_{w\sim v,\ w\in C}B_w.
\]
Consequently
\[
 D_C\le\E\!\left[\1_C\min\left\{1,
                  \left(\frac{2\Delta P}{kP}\right)_+\right\}\right].
\]
An ordinary norm bound on $\Delta P$ does not automatically control
this relative quantity. Extending the theorem to arbitrary small-degree
heads, and transferring an appropriate extension to the actual adaptive
posterior, remain separate open steps.

\section{Weighted homogeneity and two-level extensions}
The argument uses a reversible nearest-neighbor walk, rather than any
special property of equal coordinate weights. Let
\[
 A_w=\sum_{i=1}^n w_i\,\mathrm{flip}_i,
 \qquad w_i>0,\qquad\sum_iw_i=1.
\]
\begin{theorem}\label{thm:weighted}
Suppose a margin-one background $H$ on $\Omega_n$ satisfies
$A_wH=(1-2a)H$, where $a\le1/5$ and
$\max_iw_i\le1/2400$. If its full-parity error is at most $17/50$,
then every admissible repair family in the model of Section~1 satisfies
\[
 \E[\1_Ee_u]\ge C_0/2+2/5.
\]
\end{theorem}
\begin{proof}
Use the raw transport $\mu_{uv}=w_iB_u$ on an error--correct edge in
coordinate $i$. Its inverse-amplitude weight is $w_i/B_v$.
After the same row normalization $s_u=\max(a_u,1/100)$ used above,
\[
 \sum_vw_{uv}^2\le100\max_iw_i\le1/24.
\]
Thus the pointwise bound \eqref{eq:pointwise}, including its treatment
of zero rows, holds unchanged. Write
$F=\E[\1_E A_w\1_C]\le R$ and
$M=\E[\1_C A_{w,E}B/B]$. Full parity conjugates $A_w$ to $-A_w$,
so on correct vertices
\[
 A_{w,E}B=(1-2a)B+A_{w,C}B.
\]
Reversibility and $x/y+y/x\ge2$ give
$\E[\1_C A_{w,C}B/B]\ge1-R-F$. Consequently
\[
 2R+M-F\ge2(1-a)(1-R)+2(R-F)\ge2(1-a)(1-R).
\]
Averaging the pointwise bound therefore gives
$\E[\1_Ee_u]\ge4T(1-a)(1-R)-T^2VR$.
The numerical estimate in the preceding section applies verbatim.
\end{proof}

There is a useful affine example, provided the signed antipodal sheet
of the kernel is retained. Let $k$ be even, $r\ge1$ odd, and
\[
 P(x)=c+Q_r(x),\qquad H(x,b)=cb+Q_r(x),
 \qquad F(x,b)=f(x)b,
\]
where $Q_r$ is homogeneous of degree $r$. Then
$H(x,1)=P(x)$ and $H(x,-1)=-P(-x)$. Since $k$ is even, the two slices
have exactly the same parity risk. Give each $x$-coordinate weight
$1/(k+r)$ and the bias coordinate weight $r/(k+r)$. Every monomial
of $H$ has weighted degree $a=r/(k+r)$.

\begin{corollary}\label{cor:affine}
Suppose $k$ is even, $r$ is odd, $r/k\le1/2399$, and
$P=c+Q_r$ has margin one and parity error at most $17/50$.
In the odd augmented cusp-cell model, its mean wrong-center repair
energy is at least $C_0/2+2/5$.
\end{corollary}
\begin{proof}
The lift $H$ has the same margin and risk, and
$a=\max_iw_i\le1/2400$. Apply Theorem~\ref{thm:weighted}.
\end{proof}

The same construction also includes two nonconstant homogeneous levels.
\begin{corollary}\label{cor:two-level}
Let $k$ be even, let $r<s$ with $r$ even and $s$ odd, and suppose
$P=P_r+P_s$ has margin one and parity error at most $17/50$.
If
\[
 \frac{s}{k+s-r}\le\frac15,
 \qquad \frac{s-r}{k+s-r}\le\frac1{2400},
\]
then the conclusion of Corollary~\ref{cor:affine} holds for $P$.
In particular, adjacent levels $r,r+1$ are covered whenever
$k+1\ge2400$ and $(r+1)/(k+1)\le1/5$.
\end{corollary}
\begin{proof}
The odd extension is $H(x,b)=bP_r(x)+P_s(x)$ and again preserves
risk because $k$ is even. Assign the ordinary coordinates weight
$1/(k+s-r)$ and the bias coordinate weight $(s-r)/(k+s-r)$.
Both homogeneous components then have weighted degree
$a=s/(k+s-r)$. The stated assumptions are exactly those of
Theorem~\ref{thm:weighted}.
\end{proof}

Here the geometry in the corollary has a precise kernel meaning.
For an odd positive definite profile $F_L$, consider
\[
 K_d(y,z)=F_L\!\left(\frac{1+y\cdot z}{d+1}\right).
\]
It is the $b=1$ slice of the augmented kernel
$\widetilde K(v,w)=F_L(v\cdot w/(d+1))$ on $\Omega_{d+1}$.
In its canonical feature space, oddness gives
$\Phi(-v)=-\Phi(v)$. The slice contains one representative from every
antipodal pair, so restriction and the odd extension
\[
 \widetilde z(x,b)=b\,z(bx)
\]
are exactly isometric. Global margin constraints are equivalent on the
two spaces. For either background above, its extension is precisely the displayed
$H$ and remains independent of all nuisance coordinates.

A bias flip is therefore a genuine distance-one edge in the augmented
kernel. In original coordinates it is an antipodal input together with
a sign reversal; the odd kernel's endpoint cusp is the same local cusp
after this reversal. The corollary applies when the cusp-cell reduction
is performed on this augmented space, with the same radial anchoring
assumption as in Section~1. It does not identify an ordinary antipodal
pair with adjacent points in a one-sheet Hamming model. Nor does it
assert that the adaptive posterior head belongs to either two-level class.

\end{document}
